
\chapter{Poisson}
\section{Simulation setup}

The Poisson stage starts from the three-dimensional SiMOS mesh, which is
converted into a QTCAD device configured for electron confinement:
\begin{lstlisting}
d = Device(mesh, conf_carriers="e")
d.set_temperature(T)
d.statistics = "FD_approx"
\end{lstlisting}
The option \texttt{conf\_carriers="e"} tells the solver that the mobile charge
density to be computed is electronic. The temperature is set to
\SI{15}{\milli\kelvin}, consistent with the cryogenic operating regime of the
double quantum dot, and the Fermi--Dirac approximation is used to describe the
occupation of the available electronic states during the self-consistent
Poisson calculation:
\[
    T = \SI{0.015}{\kelvin} = \SI{15}{\milli\kelvin}.
\]

The electrostatic potential is controlled by the gate surfaces defined in the
geometry rather than by assigning arbitrary voltages inside the device volume.
Each surface declared with \texttt{new\_gate\_bnd} creates a boundary condition
for one metallic gate, and all gates share the same work-function offset,
\[
    \Phi_m = \chi_{\mathrm{Si}} + \frac{E_{g,\mathrm{Si}}}{2},
\]
corresponding to a mid-gap silicon reference. The different voltages in
Table~\ref{tab:poisson_voltages} then represent the intentional tuning knobs of
the device: reservoir gates accumulate carriers, barrier gates control tunnel
coupling, and plunger gates tune the energy levels of the two dots.

\begin{table}[H]
\centering
\begin{tabular}{|>{\raggedright\arraybackslash}p{5.5cm}|c|>{\raggedright\arraybackslash}p{6.1cm}|}
\hline
Boundary & Voltage (V) & Role \\
\hline
\texttt{surface\_Y1\_gate} & 0.535 & Longitudinal confinement gate. \\
\hline
\texttt{surface\_Y2\_gate} & 0.535 & Longitudinal confinement gate, biased symmetrically with Y1. \\
\hline
\texttt{surface\_Reservoir\_L\_X\_gate} & 1.300 & Left reservoir accumulation. \\
\hline
\texttt{surface\_Reservoir\_R\_X\_gate} & 1.300 & Right reservoir accumulation. \\
\hline
\texttt{surface\_Barrier\_L\_X\_gate} & 0.040 & Left tunnel barrier. \\
\hline
\texttt{surface\_Barrier\_M\_X\_gate} & 0.400 & Inter-dot barrier. \\
\hline
\texttt{surface\_Barrier\_R\_X\_gate} & 0.040 & Right tunnel barrier. \\
\hline
\texttt{surface\_Plunger\_L\_X\_gate} & 0.700 & Left dot plunger. \\
\hline
\texttt{surface\_Plunger\_R\_X\_gate} & 0.7001990578 & Right dot plunger, detuned by about \SI{199}{\micro\volt}. \\
\hline
\end{tabular}
\caption{Gate boundary voltages used in the Poisson simulation.}
\label{tab:poisson_voltages}
\end{table}

Therefore, the simulation does not fix the carrier population through explicit
source or drain ohmic contacts. The reservoirs are still present in the mesh
and in the material definition, but their role is imposed indirectly through
the high reservoir-gate voltages. This means that the formation of the
accumulated regions and their coupling to the dot are determined by the
electrostatic solution rather than by a separate ohmic-contact boundary
condition.

\section{Materials and regions}

Before the boundary conditions are applied, each mesh region is assigned a
material. This assignment fixes the dielectric response used by the Poisson
equation; in silicon it also provides the band parameters needed to compute the
electron density. The three QTCAD material objects used in the device are
reported in Table~\ref{tab:poisson_materials}.

\begin{table}[H]
\centering
\begin{tabular}{|>{\raggedright\arraybackslash}p{3.0cm}|>{\raggedright\arraybackslash}p{3.0cm}|>{\raggedright\arraybackslash}p{8.0cm}|}
\hline
QTCAD material & Physical material & Use in the device \\
\hline
\texttt{mt.Si} & Silicon & Semiconductor body, reservoirs, and quantum regions. Some silicon regions are n-doped with \(N_D=5\times10^{24}\,\mathrm{m^{-3}}\), corresponding to \(5\times10^{18}\,\mathrm{cm^{-3}}\); the p-doping is set to zero. \\
\hline
\texttt{mt.SiO2} & Silicon dioxide & Oxide regions immediately above the quantum layer and in the upper central stack. These regions act as insulating dielectrics and do not contribute mobile carriers. \\
\hline
\texttt{mt.Al2O3} & Aluminum oxide & Top oxide and gate-stack regions. The electrostatic action of the gates is imposed separately through the gate boundary conditions. \\
\hline
\end{tabular}
\caption{Materials appearing in \texttt{SiMOS\_poisson.py}.}
\label{tab:poisson_materials}
\end{table}

The same region labels are then used to identify the central quantum-dot stack:
\[
\texttt{dot\_region} =
\{\texttt{region\_pre\_quantum\_B\_mid},
\texttt{region\_quantum\_B\_mid},
\texttt{region\_post\_quantum\_B\_mid}\}.
\]

The definition of \texttt{dot\_region} is therefore a selection of the central
quantum-dot volume, not an additional material assignment. By excluding this
region from the classical charge-density calculation, QTCAD avoids treating
the dot as a continuous electron gas with strong screening. This is important
because the central region is expected to host only a small number of confined
electrons, so its electrostatic behavior must remain compatible with a later
quantum treatment instead of being fully classical.

\section{Poisson solver parameters}

The solver is instantiated with a geometry file,
\begin{lstlisting}
p = PoissonSolver(d, solver_params=params_poisson, geo_file=path_geo)
\end{lstlisting}
so QTCAD uses the adaptive-mesh non-linear Poisson solver. The meaning of each
parameter in the script is summarized in Table~\ref{tab:poisson_params}.

\begin{table}[H]
\centering
\begin{tabular}{|>{\raggedright\arraybackslash}p{3.2cm}|>{\raggedright\arraybackslash}p{2.2cm}|>{\raggedright\arraybackslash}p{8.5cm}|}
\hline
Parameter & Value & Meaning in this simulation \\
\hline
\texttt{tol} & \(10^{-7}\) & Convergence tolerance for the self-consistent Poisson loop. With the default absolute error criterion, the iteration stops when the maximum change in electrostatic potential between successive iterations is below \(10^{-7}\,\mathrm{V}\). \\
\hline
\texttt{initial\_ref\_factor} & 0.55 & Initial adaptive-refinement factor. It is dimensionless; values closer to zero force stronger mesh refinement after non-convergence, while values closer to one refine more conservatively. \\
\hline
\texttt{final\_ref\_factor} & 0.85 & Refinement factor used once the solver error has saturated. The high value means the final adaptive refinements are intentionally mild, avoiding an unnecessarily large mesh. \\
\hline
\texttt{min\_nodes} & 0 & No lower bound is imposed on the total number of mesh nodes. If the electrostatic solution converges, the adaptive procedure is allowed to stop even for a relatively small mesh. \\
\hline
\texttt{maxiter} & 1000 & Maximum number of self-consistent non-linear Poisson iterations before the solve is considered unsuccessful. The high value leaves room for convergence at the low operating temperature used in the simulation. \\
\hline
\texttt{refined\_region} & \texttt{dot\_region} & Region in which QTCAD enforces a maximum local mesh characteristic length. Here the enforced refinement is focused on the central dot stack. \\
\hline
\texttt{h\_refined} & 0.3 & Maximum characteristic mesh length in the regions listed in \texttt{refined\_region}. Since the mesh is loaded with a \(10^{-9}\) scaling factor and the Gmsh geometry is in nanometers, this corresponds to \(\SI{0.3}{\nano\meter}\). \\
\hline
\texttt{max\_nodes} & \(10^{10}\) & Upper bound on the number of mesh nodes allowed during adaptive refinement. This value is effectively non-limiting for the present calculation; the practical bound will be memory and runtime. \\
\hline
\end{tabular}
\caption{Meaning of the Poisson solver parameters used in the adaptive calculation.}
\label{tab:poisson_params}
\end{table}

Taken together, these settings make the solution accurate where the confinement
potential is most sensitive, without forcing the whole mesh to the same
resolution. The small tolerance controls the self-consistent potential, while
\texttt{h\_refined} concentrates the finest elements in \texttt{dot\_region}.
Away from the central stack, the relatively mild refinement factors allow the
adaptive solver to keep a coarser mesh where the potential varies more smoothly.
